\section{Data}
\label{sec:data}

\policyengine{} UK's microsimulation runs on an enhanced version of the
Family Resources Survey (\frs{}), constructed and calibrated by
\populace{}, \policyengine{}'s open-source population data stack. This
section documents the base survey, the enhancement pipeline, and the
calibration procedure, citing the \populace{} build code that implements
each step. \todo{Confirm this section's framing is consistent with however
\populace{}'s own model/data paper (if drafted concurrently, per the
imputation-paper and populace-be-microsim threads in the wider publication
program) describes the same pipeline, to avoid the two papers diverging on
the same facts.}

\subsection{Base survey: the Family Resources Survey}

The \frs{} is \dwp{}'s annual household survey of income, benefit receipt,
and demographic characteristics, produced for HBAI and other official
statistics. \todo{Standard FRS citation (DWP/ONS or UK Data Service), sample
size, and years currently supported by policyengine-uk -- confirm supported
years from policyengine\_uk/programs.yaml's verified\_start\_year fields
(earliest is 2013, for council\_tax\_reduction; 2019 is the earliest across
the bulk of programs, but is not the single earliest value in the file as
of 2026-07-04) and any dataset-loading code, not from memory.}

\subsection{Enhancement: SPI-informed synthetic high-income records}
\label{sec:data-efrs}

The \frs{} undersamples high incomes relative to \hmrc{}'s Survey of
Personal Incomes (\spi{}), a limitation shared by household surveys
generally. \populace{}'s UK build pipeline addresses this with an
``enhanced FRS'' (\efrs{}) construction documented in
\texttt{populace/packages/populace-build/src/populace/build/uk\_runtime/spi\_support.py}:
a zero-weight copy of the \frs{} is created and filled with \spi{}-trained
quantile-regression-forest income imputations (employment income,
self-employment income, savings interest, dividends, private pension
income, property income, plus Gift Aid and qualifying investment gifts),
and calibration (Section~\ref{sec:data-calibration}) is left free to
upweight these synthetic high-income rows where doing so improves fit to
\spi{} totals. A second imputation stage fills FRS-only fields (pension
contributions, tax-free savings income, Universal Credit receipt, and
others) on the synthetic rows so they do not inherit an FRS donor's
unrelated benefit or pension profile. \todo{Cite the specific QRF
specification (predictor set, training target) once confirmed against the
\texttt{populace-fit} code the sibling imputation paper documents -- this
section should point to that paper's methodology rather than re-derive it,
per PolicyEngine's stated preference to keep this description a summary of
what is proven elsewhere.}

The individual-level accuracy of this SPI-informed approach against SPI-
reported income tax liabilities is validated directly in
Section~\ref{sec:validation-spi}.

\subsection{Calibration}
\label{sec:data-calibration}

\policyengine{} UK's data is calibrated -- reweighted to match a set of
administrative and forecast targets -- via \populace{}'s calibration
package (\texttt{populace.calibrate}, formerly a
\texttt{policyengine\_uk.data.datasets.frs.calibration.loss} module
referenced in the model's validation notebook; that historical import path
had already been superseded by the \populace{} migration as of this
writing, and this section should cite the current location precisely once a
canonical target list is agreed). Calibration targets are defined in
\texttt{populace/packages/populace-build/src/populace/build/uk\_runtime/local\_targets.py}
and include, at minimum, the aggregate income tax liability implied by
official \obr{} costings (the loss module referenced program-level budgetary
impact parameters such as
\texttt{programs.income\_tax.budgetary\_impact.by\_country.UNITED\_KINGDOM}).
\todo{Enumerate the full current target list (which programs, which
geography levels -- the local\_targets.py module also supports local-area
row-wise geography per the UK local-area build tests) and the calibration
method (the populace README describes \texttt{populace-calibrate} as
supporting APG and L0-regularized approaches; confirm which the UK build
uses in production).}

\subsection{Nowcasting and uprating}

Direct input variables (employment income, benefit receipt, rents, council
tax, etc.) are uprated from the survey's collection year to the simulation
year using \obr{} economic indices, documented in the model's HBAI mapping
page (\texttt{docs/book/validation/hbai.md}), which tabulates the specific
uprating index used for every income, subtraction, and housing-cost
component in the \hbai{} income definition (e.g. \texttt{employment\_income}
uprated on \texttt{gov.obr.per\_capita.employment\_income}; rents uprated on
outturn regional data through 2025 and \obr{} aggregate rent forecasts
thereafter). \policyengine{}'s own growth-rate assumptions are documented
separately in \texttt{docs/book/assumptions/growthfactors.md}.
\texttt{docs/book/assumptions/rf-nowcasting-methodology.md} is a distinct,
paired reference: it summarizes the Resolution Foundation's own (external,
third-party) nowcasting methodology -- not \policyengine{}'s -- for
comparison, with the high-level differences between the two approaches
tabulated in \texttt{docs/book/assumptions/nowcasting-comparison.md}.
\todo{Summarize \texttt{growthfactors.md} and, if a citable comparison is in
scope for this paper, the specific points of difference from RF's approach
in \texttt{nowcasting-comparison.md} -- neither has been read in full for
this outline; only their existence and stated purpose (confirmed via
\texttt{gh api} directory listing and reading \texttt{rf-nowcasting-methodology.md}
in full) are established here.}

\subsection{Release input-coverage contract}
\label{sec:data-coverage}

Every input column the UK national build exports to the engine is governed
by an enforced release contract, merged as PolicyEngine/populace\#420
(``UK release input-coverage contract: 145-required gate, effective-mass
semantics, and the real-donor HMRC replay,'' merged 2026-07-13): the
current contract requires \textbf{145} inputs and carries \textbf{0}
reviewed exclusions. Coverage is defined on \emph{effective mass}, not raw
nonzero counts: a required input must carry weighted signal on at least
one part per million of its owning entity's effective population mass, so
a column populated only on zero-weight rows does not count as covered. The
gate fails closed: a release candidate that cannot populate a required
input fails the release, and the only alternative status is a reviewed,
evidence-backed exclusion -- of which the current contract carries none.
Degenerate-required, wrong-entity, dormant/stale-exclusion, and
missing-column conditions each fail the gate with named columns (per the
merged PR record and the adjudication register committed at
\texttt{UK\_COVERAGE\_PROGRESS.md} in the \populace{} repository).

Exclusions themselves are monitored rather than grandfathered. The
register opened at 143 required inputs and 2 reviewed exclusions, because
two charitable-giving columns (\texttt{gift\_aid} and
\texttt{charitable\_investment\_gifts}) carried signal only on zero-weight
synthetic rows; when a rebuilt \spi{} stage gave both columns weighted
signal above the floor, the stale-exclusion check itself demanded their
promotion to hard requirements, producing the current 145/0 state
(\texttt{UK\_COVERAGE\_PROGRESS.md}). For scale, the equivalent US
register in the same monorepo (\texttt{COVERAGE\_PROGRESS.md}, verified
2026-07-15) stands at 158 required inputs and 8 reviewed exclusions; the
UK contract carries zero.

\subsection{Local-area geography}
\label{sec:data-geography}

\populace{}'s UK build adds a local-geography spine mirroring its US
census-block ladder: anchor every household at the finest published census
geography, derive every other layer deterministically from that anchor,
and keep one national dataset that is filterable at any grain rather than
publishing per-area files. The design is recorded in
PolicyEngine/populace\#349 and the implementation is merged as
populace\#354 (``UK output-area-anchored geography ladder (US
block-ladder pattern),'' merged 2026-07-09):
\texttt{populace/packages/populace-build/src/populace/build/uk\_runtime/geography\_ladder.py}
anchors each household at a 2021 Census output area via two seeded draws
-- a Westminster parliamentary constituency (2024 boundaries) sampled
within the household's calibrated region proportional to constituency
household counts, then an output area within that constituency
proportional to census population -- and derives every other layer (LSOA
and MSOA by structural nesting; local authority district, ward, and ITL
levels by pinned ONS lookups) from the anchor with no independent
randomness.

The merged module carries the discipline the US ladder established:
per-layer source vintages are recorded in the ladder artifact and are
mandatory (the loader refuses an artifact missing any, and assignment
refuses a ladder whose constituency vintage differs from the one the
constituencies were assigned under); exported columns are engine inputs or
plain data columns, never persisted formula outputs; and a
release-blocking gate checks GSS-code structure, ITL-prefix consistency,
and a London-mass floor. At merge the ladder covers England and Wales
end-to-end; a household whose region is absent from the ladder is refused
by name rather than silently part-joined, with Scotland and Northern
Ireland rungs to follow once their lookup vintages are pinned (per the PR
record). \todo{Populace\#354 merges the ladder runtime, artifact builder,
and gates, but deliberately defers wiring the assignment into the UK build
behind the shared geography-spine schema to a follow-up -- the released UK
artifact does not yet carry ladder-assigned output areas. Cite the
follow-up PR and update this subsection once the assignment is wired into
a released artifact.}

\subsection{Data versioning and promotion}
\label{sec:data-versioning}

\todo{State the currently certified \populace{} bundle version for the UK
build (per internal tracking, certified defaults were "4.18.7 (dense)" for
both US and UK as of early July 2026 -- verify this is still current and
cite the populace release/tag directly rather than an internal memory note,
since this paper must cite public, checkable artifacts).}

The UK default dataset was promoted through an adjudicated comparison
rather than by recency. The current default, \texttt{populace\_uk\_2023},
was certified in PolicyEngine/\texttt{policyengine.py}\#427 (``Certify UK
Populace default dataset,'' merged 2026-06-19) on a matched-$N$ holdout
comparison against the incumbent enhanced \frs{}: per the PR record,
holdout loss 0.1239 vs.\ 0.3784, per-target wins 79 to 70, mean absolute
relative target error 1.49\% with the worst target miss at 23.1\% against
a 25\% release gate, and the incumbent's prior dataset names retained as
resolvable aliases. A proposal to formalize this promotion gate as a
pre-registered certification suite run at every default flip, together
with a recorded retroactive adjudication of the UK promotion on broader
metrics, is \texttt{policyengine.py}\#462 (open at this writing) and is
summarized in Section~\ref{sec:validation-promotion}.

\subsection{Licensing and data handling}
\label{sec:data-licensing}

The source microdata the pipeline consumes -- the \frs{} tables and the
\spi{} public use tape -- are licensed through the UK Data Service and are
handled outside version control: licensed row-level files never enter the
\populace{} repository or its git history, tracked build artifacts are
aggregate-only (the committed replay and gate reports contain no row-level
donor data), and staging datasets remain untracked
(PolicyEngine/populace\#420; \texttt{UK\_COVERAGE\_PROGRESS.md}). Build
records instead pin each licensed input by checksum, so provenance is
verifiable without redistributing the data.
